\documentclass[10pt,a4paper]{amsart}
\usepackage[margin=19mm]{geometry}
\usepackage{amsmath,amssymb,mathtools}
\usepackage[scr=rsfs,cal=euler]{mathalfa}
\usepackage{xfrac}
\usepackage[unicode,hidelinks,bookmarks=false]{hyperref}
\newcommand{\ie}{i.e.\ }
\newcommand{\cf}{cf.\ }
\newcommand{\resp}{respectively\ }
\newcommand{\Subsection}{\subsection}
\providecommand{\nobreakdash}{\nobreak\hbox{-}}
\setlength{\emergencystretch}{2em}
\multlinegap0pt
\usepackage{bm}
\usepackage{bbm}
\usepackage{dsfont}
\usepackage{xspace}
\usepackage{cancel}
\usepackage{mathtools}
\usepackage{amsthm}
\makeatletter
\@ifundefined{lemma}{\newtheorem{lemma}{Lemma}}{}
\@ifundefined{theorem}{\newtheorem{theorem}{Theorem}}{}
\makeatother
\title[Coercivity and Local Convergence of Physical Learning in Lin]{Coercivity and Local Convergence of Physical Learning in Linear Circuits}
\author{Joshua A. McGinnis, Xinbo Li, Yoichiro Mori}
\date{Source: arXiv:2606.15443v1 (CC BY 4.0)}
\begin{document}
\maketitle

\noindent\emph{Source and licence.} Coercivity and Local Convergence of Physical Learning in Linear Circuits. Version arXiv:2606.15443v1 (CC BY 4.0), distributed under the Creative Commons Attribution 4.0 International licence, as recorded in the arXiv licence metadata for this exact version and shown on the abstract page https://arxiv.org/abs/2606.15443v1. Full rights notice: licenses/arxiv-2606.15443-CC-BY-4.0.txt.\par\smallskip

\noindent\textit{Editorial scope.} Counted unit: the Theorem whose source label is \texttt{thm:local\_convergence} in the article (source0.tex lines 922--924, source label \texttt{thm:local\_convergence}), delivered together with the article's own complete original proof of it (source0.tex lines 926--1022, one \texttt{proof} environment of 97 lines). No mathematics is written, repaired or re-derived: statement and proof are the article's own text line for line. Required context retained uncounted: thm:EP\_coercivity (lines 660--667); lem:CL\_coercivity\_iff (lines 692--699); thm:CL\_coercivity (lines 726--737). Every definition, display and label of those lines is retained unchanged. Omitted from this package and not redistributed: 1--659, 668--691, 700--725, 738--921, 925--925, 1023--1232. Nothing inside a retained excerpt is omitted or altered: every retained body is delivered exactly as the article's lines are written, so no omission receipt of any kind accompanies this package. Citations: the retained excerpts carry no citation command at all, so no bibliography entry of the article is transcribed and no reference is redistributed. Reference adaptation: none was needed and none was made; every reference inside the delivered unit resolves inside it, so no unresolved reference or citation is produced. Printed slips kept literal: no printed word, symbol, hypothesis or number of the counted statement or of its proof was altered. Layout and class: the article's own class is \texttt{siamart} with packages including makecell, mathrsfs, circuitikz, amssymb, colortbl; the wrapper is the lane's pinned \texttt{amsart} editorial wrapper at 10pt on A4, which restates the theorem-like environments the retained text needs and adds the article's own macro definitions that the retained text needs (none), with extra wrapper packages (bm, bbm, dsfont, xspace, cancel, mathtools). The delivered numbering is the unit's own. Assets: no retained span contains a figure environment, a graphics-inclusion command, a PDF-page inclusion, a table or a TikZ picture; no image file is copied into this package, the package declares no asset dependency and no image file is redistributed with it. Licence: version arXiv:2606.15443v1 of this article is distributed under the Creative Commons Attribution 4.0 International licence, as recorded in the arXiv licence metadata for this exact version and shown on the abstract page https://arxiv.org/abs/2606.15443v1. A full rights notice is delivered at licenses/arxiv-2606.15443-CC-BY-4.0.txt. Characters: every retained excerpt is pure ASCII, so the delivered engine, XeLaTeX with its default Latin Modern text font, reports no missing character.
\begin{theorem}[Coercivity of EP]
\label{thm:EP_coercivity}
    Suppose $k_0>0$ and that $\mathscr{D}'^\top\hat{L}^{-1}\hat{Q}$ has full column rank at $k_0$. Set $\lambda := \lambda_{\min}(\Lambda(k_0))$ and $c := \sigma_{\min}\bigl(\mathscr{D}'^\top\hat{L}^{-1}\hat{Q}\bigr)\big|_{k_0}$. Then there is a neighborhood $U$ of $k_0$ such that for all $k \in U$ the loss $\Phi = \frac{1}{2}\|r\|^2$ satisfies
\begin{equation}
    \dot{\Phi} \leq -2\lambda c^2\,\Phi.
\end{equation}
In particular, $\Phi$ decays exponentially as long as $k$ remains in $U$.
\end{theorem}

\begin{lemma}
\label{lem:CL_coercivity_iff}
Suppose $k_0>0$. Then there exists a positive constant $c_D$ such that
\[
\sum_{i=1}^M\big\langle R Q D\tilde{r},\partial_{k_i}\nabla_x G\big\rangle^2 \geq c_D\,\|\tilde{r}\|_2^2
\]
holds for all $\tilde{r} \in \mathbb{R}^O$ and all $k$ in a neighborhood of $k_0$ if and only if $\mathscr{D}'^\top\hat{L}^{-1}\hat{Q}$ has full column rank at $k_0$.
\end{lemma}

\begin{theorem}[Coercivity of CL]
\label{thm:CL_coercivity}
    Suppose $k_0>0$ and that $\mathscr{D}'^\top\hat{L}^{-1}\hat{Q}$ has full column rank at $k_0$. Let $\lambda_D, c_D > 0$ denote the coercivity constants for $\Phi_D$ at $k_0$, i.e., constants such that $\|\nabla\Phi_D\|^2 \ge 2\lambda_D c_D^2\,\Phi_D$ on a neighborhood of $k_0$ (these are assembled from Lemma~\ref{lem:CL_coercivity_iff} and the bounds on $\lambda_{\min}(D), \lambda_{\max}(D)$). Then there is a neighborhood $U$ of $k_0$ and a constant $C_D > 0$ such that for all $k \in U$ the weighted loss $\Phi_D = \frac{1}{2}\langle r, Dr\rangle$ satisfies
\begin{equation}
    \dot{\Phi}_D  \leq -2\lambda_D c_D^2\,\Phi_D + C_D\,\Phi_D^{3/2}.
\end{equation}
In particular, if $k_0 \in \mathcal{S}$ (so that $\Phi_D(k_0) = 0$), then by shrinking $U$ if necessary,
\begin{equation}
    \dot{\Phi}_D \leq -\lambda_D c_D^2\,\Phi_D
\end{equation}
for all $k \in U$.
\end{theorem}

\begin{theorem}\label{thm:local_convergence}
Suppose that $k^*$ is a point on the solution manifold $\mathcal{S}:=\{k \in \mathbb{R}^M \mid r(k) = 0\}$ such that $\nabla_k r$ has full row rank. (For EP and CL, this holds exactly when the active-incidence conditions of Theorems~\ref{thm:EP_coercivity} and~\ref{thm:CL_coercivity} are satisfied.) Then there exists a positive real number $\delta$ such that for all $\|k(0) - k^*\| < \delta$, we have $\lim_{t \to \infty} k(t) \in \mathcal{S}$ and $\|r(k(t))\|$ decays exponentially.
\end{theorem}

\begin{proof}
Since $\nabla_k r(k^*)$ has full row rank, we may reorder coordinates and write
\[
k = (u,v), \qquad k^* = (u^*, v^*),
\]
with $u \in \mathbb{R}^{M-O}$ and $v \in \mathbb{R}^{O}$, such that
\[
E := \frac{\partial r}{\partial v}(k^*)
\]
is invertible.

By the Implicit Function Theorem, there exist open neighborhoods $U \subset \mathbb{R}^{M-O}$ of $u^*$ and $V \subset \mathbb{R}^{O}$ of $v^*$, and a $C^1$ map
\[
\psi : U \to V
\]
such that
\[
\mathcal{S} \cap (U \times V) = \{ (u, \psi(u)) : u \in U \}.
\]
In particular, $U \times V$ is an open neighborhood of $k^*$.

By shrinking $U \times V$ if necessary, we may assume the following hold uniformly for all $k \in U \times V$:
\begin{enumerate}
    \item[(i)] $\dfrac{\partial r}{\partial v}(k)$ is invertible with $\left\|\left(\dfrac{\partial r}{\partial v}(k)\right)^{-1}\right\| \le C_1$,
    \item[(ii)] $|\dot{k}| \le C_2 \|r(k)\|$,
    \item[(iii)] For EP: $\dot{\Phi} \leq -2\lambda c^2\,\Phi$ (Theorem~\ref{thm:EP_coercivity}). For CL: $\dot{\Phi}_D \leq -\lambda_D c_D^2\,\Phi_D$ (Theorem~\ref{thm:CL_coercivity}), with $D$ having uniformly bounded eigenvalues so that $\tfrac{1}{2}\lambda_{\min}(D)\|r\|^2 \leq \Phi_D \leq \tfrac{1}{2}\lambda_{\max}(D)\|r\|^2$,
    \item[(iv)] $\psi(U) \subset V$.
\end{enumerate}


Note that for $k = (u,v) \in U \times V$,
\begin{equation}
\operatorname{dist}(k,\mathcal S)
\le \|(u,v)-(u,\psi(u))\|
= \|v-\psi(u)\|.
\label{eq:dist-bound-v}
\end{equation}

Using the Fundamental Theorem of Calculus along the segment from $(u,\psi(u))$ to $(u,v)$,
\[
r(u,v)=r(u,v) - r(u,\psi(u))
= 
\int_0^1
\frac{\partial r}{\partial v}\big(u, \psi(u) + t(v-\psi(u))\big)\, dt \,(v-\psi(u)).
\]
Define
\[
A :=
\int_0^1
\frac{\partial r}{\partial v}\big(u, \psi(u) + t(v-\psi(u))\big)\, dt.
\]
Then
\begin{equation}
r(u,v) = A (v-\psi(u)).
\label{eq:r-factored}
\end{equation}
By (i), $A$ is a small perturbation of $E$, and hence invertible with $\|A^{-1}\| \le 2C_1$. From \eqref{eq:r-factored},
\[
\|v-\psi(u)\| \le \|A^{-1}\| \|r(u,v)\| \le 2C_1 \|r(k)\|.
\]

Combining with \eqref{eq:dist-bound-v},
\begin{equation}
\operatorname{dist}(k,\mathcal S) \le 2C_1 \|r(k)\|
\quad \text{for all } k \in U \times V.
\label{eq:dist-bound-r}
\end{equation}

As long as $k(t) \in U \times V$, conditions (i)--(iv) hold. By (iii), coercivity gives
\[
\dfrac{d}{dt}\|r(k(t))\|^2 \le -2\lambda c^2\, \|r(k(t))\|^2,
\]
and hence
\[
\|r(k(t))\| \le e^{-\lambda c^2\, t}\, \|r(k(0))\|
\]
for EP. For CL, the analogous inequality has
\[\|r(k(t))\| \le \sqrt{\tfrac{\lambda_{\max}(D)}{\lambda_{\min}(D)}}\,e^{-\lambda_D c_D^2\, t/2}\, \|r(k(0))\|.\]
In either case, there exist $K\ge 1$ and $\nu>0$ such that \[\|r(k(t))\| \le Ke^{-\nu t}\|r(k(0))\|\] (with $K=1,\,\nu=\lambda c^2$ for EP; $K=\sqrt{\lambda_{\max}(D)/\lambda_{\min}(D)},\,\nu = \lambda_D c_D^2/2$ for CL).          

Since $r(k^*) = 0$ and $r$ is $C^1$, there exists $C_3 > 0$ such that $\|r(k)\| \leq C_3\|k - k^*\|$ on $U \times V$. By (ii), the total displacement of $k$ is controlled:
\[
\|k(t) - k(0)\| \le C_2 \int_0^t \|r(k(s))\|\, ds \le \dfrac{KC_2}{\nu}\, \|r(k(0))\| \leq \dfrac{KC_2 C_3}{\nu}\, \delta.
\]
By the triangle inequality,
\[
\|k(t) - k^*\| \leq \|k(t) - k(0)\| + \|k(0) - k^*\| < \left(\dfrac{KC_2 C_3}{\nu}+1\right)\delta.
\]
Choose $\delta$ small enough that $\left(\dfrac{KC_2 C_3}{\nu}+1\right)\delta < \operatorname{dist}(k^*, \partial(U \times V))$. By continuity, $k(t) \in U \times V$ for $t$ in some maximal interval $[0,T)$. The bounds above hold on $[0,T)$ and show $k(t)$ remains bounded away from $\partial(U \times V)$, so $T = \infty$.

It remains to show $\lim_{t \to \infty} k(t)$ exists. For $0 \le s \le t$, the same estimate as above gives
\[
\|k(t) - k(s)\| \le C_2 \int_s^t \|r(k(\tau))\|\, d\tau
\le \dfrac{KC_2}{\nu}\, e^{-\nu\, s}\, \|r(k(0))\|,
\]
which tends to $0$ as $s \to \infty$ uniformly in $t \ge s$. Hence $\{k(t)\}$ is Cauchy and $k_\infty := \lim_{t \to \infty} k(t)$ exists. Since $\|r(k(t))\|$ decays exponentially for all $t \ge 0$, \eqref{eq:dist-bound-r} gives $\operatorname{dist}(k(t),\mathcal S) \to 0$, and since $\mathcal{S}$ is closed, $k_\infty \in \mathcal{S}$.
\end{proof}

\end{document}
